% Part: lambda-calculus
% Chapter: church-rosser
% Section: beta-reduction

\documentclass[../../../include/open-logic-section]{subfiles}

\begin{document}

\olfileid[fr]{lam}{cr}{b}
\olsection{\texorpdfstring{$\beta$}{β}-réduction}

\begin{lem}\ollabel{lem:one-par}
Si $M \bredone M'$, alors $M \bredpar M'$.
\end{lem}

\begin{proof}
Raisonnons sur la dérivation de la
contraction~$M\bredone M'$. Si le rédex
est le terme entier, écrivons
$M=(\lambd[x][N])Q$ et
$M'=\Subst{N}{Q}{x}$. La réflexivité
$N\bredpar N$ et $Q\bredpar Q$ vient du
\olref[pb]{thm:refl}; la règle
\olref[pb]{defn:bredpar4} donne donc
$M\bredpar M'$. Si la contraction a lieu
dans le corps d'une abstraction, appliquons
l'hypothèse d'induction et la règle
\olref[pb]{defn:bredpar2}. Si elle a lieu
dans l'un des deux membres d'une application,
utilisons l'hypothèse d'induction sur ce
membre, la réflexivité sur l'autre et la règle
\olref[pb]{defn:bredpar3}.
\end{proof}

\begin{lem}\ollabel{lem:par-red}
Si $M \bredpar M'$, alors $M \bred M'$.
\end{lem}

\begin{proof}
Raisonnons par induction sur la dérivation
de $M\bredpar M'$.
\begin{enumerate}
\item Si la dernière étape applique la règle
  \olref[pb]{defn:bredpar1}, alors
  $M=M'=x$ et $x\bred x$ par réflexivité.
\item Si la dernière étape applique la règle
  \olref[pb]{defn:bredpar2}, alors
  $M=\lambd[x][N]$ et
  $M'=\lambd[x][N']$, avec
  $N\bredpar N'$. Par hypothèse
  d'induction, $N\bred N'$. La
  compatibilité de chaque contraction
  avec l'abstraction donne
  $\lambd[x][N]\bred\lambd[x][N']$.
\item Si la dernière étape applique la règle
  \olref[pb]{defn:bredpar3}, alors
  $M=PQ$ et $M'=P'Q'$, avec
  $P\bredpar P'$ et $Q\bredpar Q'$.
  Les hypothèses d'induction donnent
  $P\bred P'$ et $Q\bred Q'$.
  Réduisons d'abord $P$ dans $PQ$,
  puis $Q$ dans $P'Q$ : nous obtenons
  $PQ\bred P'Q'$.
\item Si la dernière étape applique la règle
  \olref[pb]{defn:bredpar4}, alors
  $M=(\lambd[x][N])Q$ et
  $M'=\Subst{N'}{Q'}{x}$, avec
  $N\bredpar N'$ et $Q\bredpar Q'$.
  Par hypothèse d'induction,
  $N\bred N'$ et $Q\bred Q'$.
  Réduisons le corps de l'abstraction,
  puis son argument, et contractons
  enfin le rédex obtenu :
  \[
    (\lambd[x][N])Q
    \bred (\lambd[x][N'])Q'
    \bredone \Subst{N'}{Q'}{x}.
  \]
  Donc $M\bred M'$.
\end{enumerate}
\end{proof}

\begin{lem}\ollabel{lem:str}
$\bred$ est la plus petite relation
transitive qui contient $\bredpar$.
\end{lem}

\begin{proof}
Soit $\xred$ la plus petite relation
transitive qui contient $\bredpar$.

Montrons d'abord que
$\bred\subseteq\xred$. Une chaîne
$M\bred M'$ de contractions~$\bredone$
se transforme, par le \olref{lem:one-par},
en une chaîne de pas~$\bredpar$; la
transitivité donne $M\xred M'$.
Si la chaîne initiale a longueur zéro,
la réflexivité de $\bredpar$, établie dans le
\olref[pb]{thm:refl}, donne aussi
$M\xred M$.

Réciproquement, montrons
$\xred\subseteq\bred$. Toute chaîne
de pas~$\bredpar$ est, par le
\olref{lem:par-red}, une chaîne de
réductions~$\bred$. La transitivité de
$\bred$ donne donc l'inclusion voulue.
\end{proof}

\begin{thm}\ollabel{thm:cr}
$\bred$ vérifie la propriété de
Church--Rosser.
\end{thm}

\begin{proof}
Par le \olref[pb]{thm:cr}, $\bredpar$ vérifie
la propriété de Church--Rosser.
Par le \olref{lem:str}, $\bred$ est sa
clôture transitive; le \olref[dap]{thm:str}
transmet donc la propriété à $\bred$.
\end{proof}

La propriété de Church--Rosser entraîne l'unicité de la forme
normale~$\beta$, lorsqu'elle existe.

\end{document}
